\section{Introduction}

The ICH M7(R2) guideline endorses two complementary in~silico approaches for assessing the mutagenic potential of pharmaceutical impurities: a rule-based expert system and a statistical QSAR model.\cite{ICH_M7_2017} This regulatory framework has spurred sustained interest in predictive models for the Ames bacterial reverse mutation assay,\cite{Ames_1973,Mortelmans_2000} with extensions to in~vitro chromosome aberration (CA) and in~vivo micronucleus (MN) tests.\cite{OECD_473,OECD_474} Reliable models could reduce animal testing and accelerate safety assessment in line with the 3Rs principle.\cite{Russell_Burch_1959}

Published Ames QSAR models report a wide spread of metrics---accuracies of 0.80--0.90 and MCC values of 0.60--0.80\cite{Hansen_2009,Xu_2012,Honma_2019,Li_2021}---but these numbers come from different splits, datasets, preprocessing pipelines, and evaluation protocols, making direct comparison difficult. At least three factors can inflate apparent performance.

\textbf{Splitting strategy.} Random splits allow structurally similar compounds (sharing the same Murcko scaffold\cite{Bemis_Murcko_1996}) to appear in both training and test sets, producing optimistic generalization estimates. Scaffold-based splitting\cite{Sheridan_2013,Wu_MoleculeNet_2018} is more realistic, but the magnitude of the random-vs-scaffold gap has not been systematically quantified across genotoxicity endpoints.

\textbf{Source mixing.} Most large Ames datasets aggregate compounds from pharmaceutical libraries and industrial chemical registries.\cite{Hansen_2009,Honma_2019} These sources differ in composition, mutagenicity prevalence, and likely in data quality and labeling conventions. When mixed, models may learn source membership rather than toxicity mechanisms---a form of source heterogeneity noted\cite{Tropsha_2010,Hanser_2016} but never systematically measured.

\textbf{Preprocessing.} Salt stripping, metal removal, and charge normalization can substantially alter molecular representations,\cite{Fourches_2010,Fourches_2016} yet these steps are often undocumented or applied uniformly across endpoints.

Most studies evaluate a single endpoint with a single split, focus on algorithm comparison,\cite{Yang_2019,Mayr_2016} rarely test cross-domain generalization, and seldom report statistical significance.\cite{Demsar_2006} The relative contributions of split method, source composition, and preprocessing to apparent performance have not been disentangled within a single consistent framework.

\textbf{Contributions.} We (i)~introduce a three-level evaluation framework---scaffold-split test, random vs.\ scaffold CV, and cross-source leave-domain-out---applied to a single locked configuration for fair comparison; (ii)~show that source heterogeneity from mixing drug-like and industrial compounds (19.5-fold prevalence gap) drove a scaffold-to-LDO MCC drop of 0.390, an order of magnitude larger than the split-method effect ($\Delta = 0.046$); (iii)~find that in-domain algorithm differences are small (MCC range = 0.063), but cross-domain robustness varies sharply (RF 0.325 vs.\ XGBoost 0.122 in the hardest direction); (iv)~report endpoint-specific preprocessing effects, with salt stripping improving the best in~vivo model by +0.22 MCC while harming in~vitro; and (v)~release an open-source pipeline with all configurations and seeds for full reproducibility.
